% Part: computability
% Chapter: recursive-functions
% Section: trees

\documentclass[../../../include/open-logic-section]{subfiles}

\begin{document}

\olfileid{cmp}{rec}{tre}
\olsection{الأشجار}

يفيد أحيانًا تمثيل الأشجار بأعداد طبيعية، مثلما يمكننا تمثيل المتتاليات
بأعداد وتمثيل خصائصها والعمليات عليها بعلائق ودوال عودية بدائية في
رموزها. وسنستعمل المتتاليات ورموزها لهذا الغرض. يمكن أن تكون الشجرة
عقدة واحدة (قد تحمل لصيقة)، أو عقدة (قد تحمل لصيقة) متصلة بعدد من
الأشجار الجزئية. وتسمى العقدة \emph{جذر} الشجرة، وتسمى الأشجار الجزئية
المتصلة بها \emph{أشجارها الجزئية المباشرة}.

نرمز الأشجار عوديًا بمتتالية $\tuple{k,d_1,\dots,d_k}$، حيث $k$ عدد
الأشجار الجزئية المباشرة، وحيث $d_1$، …،~$d_k$ رموز هذه الأشجار
الجزئية. وإذا حملت العقد لصائق، أمكن إدراجها بعد الأشجار الجزئية
المباشرة. وهكذا نرمز إلى شجرة لا تتألف إلا من عقدة واحدة ذات لصيقة~$l$
بـ$\tuple{0,l}$، ونرمز إلى شجرة تتألف من جذر (ذو لصيقة~$l_1$) متصل
بعقدتين منفردتين (ذواتي اللصيقتين $l_2$ و$l_3$) بـ
$\tuple{2,\tuple{0,l_2},\tuple{0,l_3},l_1}$.

\begin{prop}
  \ollabel{prop:subtreeseq}
  الدالة $\fn{SubtreeSeq}(t)$، التي تعيد رمز متتالية عناصرها رموز جميع
  الأشجار الجزئية للشجرة ذات الرمز~$t$، عودية بدائية.
\end{prop}

\begin{proof}
  لاحظ أولًا أن $\fn{ISubtrees}(t)=\fn{subseq}(t,1,(t)_0)$ عودية
  بدائية، وأنها تعيد رموز الأشجار الجزئية المباشرة لشجرة~$t$. ويمكننا
  الآن تعريف دالة مساعدة $\fn{hSubtreeSeq}(t,n)$ تحسب متتالية جميع
  الأشجار الجزئية التي لا يزيد بعدها عن الجذر على $n$ عقد. ومتتالية الأشجار الجزئية
  لـ~$t$ التي تبعد $0$ عقد عن الجذر، أي التي تبدأ عند جذر~$t$، هي
  المتتالية التي لا تتألف إلا من~$t$. وللحصول على متتالية جميع الأشجار
  الجزئية حتى المستوى~$n+1$ للشجرة $t$، نصل الأشجار الجزئية حتى المستوى
  $n$ بمتتالية تتألف من جميع الأشجار الجزئية المباشرة لأشجار المستويات
  حتى $n$. وللحصول على قائمة بها كلها، لاحظ أنه إذا كانت $f(x)$ دالة عودية
  بدائية تعيد رموز متتاليات، فإن الدالة $g_f(s,k)$ التي تصل قيم $f$
  لأول $k$ من عناصر $s$ عودية بدائية أيضًا:
    \begin{align*}
      g_f(s, 0) & = \emptyseq\\
      g_f(s, k+1) & = g_f(s, k) \concat f((s)_k)
    \end{align*}
    فمثلًا، إذا كانت $s$ متتالية أشجار، فإن
    $h(s)=g_{\fn{ISubtrees}}(s,\len{s})$ تعطي متتالية الأشجار الجزئية
    المباشرة لعناصر~$s$. ويمكننا استعمالها لتعريف
    $\fn{hSubtreeSeq}$ بـ
    \begin{align*}
      \fn{hSubtreeSeq}(t, 0) & = \tuple{t} \\
      \fn{hSubtreeSeq}(t, n+1) & = \fn{hSubtreeSeq}(t, n) \concat
      h(\fn{hSubtreeSeq}(t, n)).
    \end{align*}
    وأقصى مستوى للأشجار الجزئية في شجرة رمزها~$t$، أي أقصى مسافة بين
    الجذر وعقدة ورقية، محدود بالرمز~$t$. ولذلك تعطي
    $\fn{hSubtreeSeq}(t,t)$ متتالية رموز جميع الأشجار الجزئية للشجرة
    ذات الرمز~$t$.
\end{proof}

\begin{prob}
  يتضمن تعريف $\fn{hSubtreeSeq}$ في برهان
  \olref[cmp][rec][tre]{prop:subtreeseq} تكرارات بوجه عام. أعط تعريفًا
  بديلًا يضمن ألا يرد رمز شجرة جزئية إلا مرة واحدة في القائمة الناتجة.
\end{prob}

\end{document}
